\scanpage{23}

se trouve de façon naturelle réduit à $\mu_n$.

On prend maintenant
\[
\begin{gathered}
C=\text{catégorie des courbes elliptiques de}\\
\text{Hasse nul sur corps alg. clos $k$ de car $p$}
\end{gathered}
\]
D'après Deuring, c'est une catégorie de Deuring \inserted{stricte sur $\mathbf Z$}, \deleted{dont} correspondant à la classe d'isomorphie des corps gauches de degré $4$ (corps de quaternions) sur $\mathbf Q$ \deleted{dont} qui \unsure{splitte} en ttes les places sauf $p$ et $\infty$ (il est donc d'ordre $2$ dans $\operatorname{Br}(\mathbf Q)$).

On tire donc, canoniquement en \deleted{\unsure{la}} termes de \deleted{\unsure{G}}\inserted{$C$} [\deleted{la donnée} \inserted{ou de} d'un tel corps sur $\mathbf Q$, une gerbe $\mathcal G$ sur $\operatorname{Spec}\mathbf Z[\frac1p]$ de \deleted{groupe} \inserted{lien} $\mu_{\unsure{2|n}}$ [$\simeq$ Gerbe $\widehat{\mathcal G}$ du \unsure{type|lien} $G_m$, $+$ représentation \inserted{inversible, $T$,} de ce dernier, de rang $2$]. On trouve de plus, si
\[
A'=\prod_{\ell\ne p,\infty}\mathbf Z_\ell,
\]
une \deleted{représentation de $\unsure{\mathcal G}_{A'}$ de rang $2$, degré $1$, soit $E$, (\unsure{donc un isom} \unsure{avec une repr.} de $\mathcal G_{A'}$ de rang $2$, degré $\unsure{1_T|1}$), avec
\[
\det E\simeq T
\]}
\inserted{banalisation de $C_{A'}$, i.e.}
\inserted{i.e. une trivialisation de $\widehat{\mathcal G}_{A'}$, $\varphi$, \deleted{\unsure{ou}}}
\margin{\fbox{i.e. foncteur fibre $\varphi$ de $\operatorname{Rep}(\widehat{\mathcal G}_{A'})$ sur $A'$}}
\inserted{muni} d'un isom
\[
\deleted{\unsure{\varphi E}}\qquad
\varphi(T)\simeq A'(1)\left(\overset{\mathrm{def}}=\prod_{\ell\ne p,\infty}\mathbf Z_\ell(1)\right)
\]
\deleted{\unsure{donc} Cela fait donc}
En d'autres termes, si $\mathcal H_{A'}$ est la gerbe des racines carrées de $A'(1)$, qui est une gerbe sous $\mu_{2,A'}$, on a un isomorphisme canonique de

\mspage{22}
